% issue #1057 的最小复现：xeCJKfntef 命令嵌套使用时的断行边界。
%
% 用 xelatex 编译。观察点：
%   （1）线型套线型 —— 内层整段无法断行，溢出右边界，日志里有 Overfull \hbox。
%   （2）线型套符号型 —— 正常断行。
%   （3）按语义分段、每段只用一个线型命令 —— 正常断行。
%
% 实测本文件只在第（1）例产生 Overfull \hbox（152.66pt；具体数值随页宽变化，
% 换用 10cm 页宽时是 276.99pt）。发布版 v3.10.4 与工作树 v3.10.5 的结果一致，
% 说明这是 ulem 的固有约束，不是某个版本的回归。
%
% 根因：ulem 只允许最外层的线型命令启动扫描。内层线型命令会复用外层已经打开的
% 扫描过程，它的正文被整段放进一个固定宽度的盒子，盒子内部不再留有断点，因此
% TeX 只能让这一段溢出。
%
% 正文一律写成字面记号，不用宏承载。写成 \CJKunderline{\BODY} 会触发另一条已知
% 限制（宏体在 ulem 扫描期间才展开），两条限制会被混为一谈。
\documentclass{article}
\usepackage[LoadFandol]{xeCJK}
\usepackage{xeCJKfntef}
\usepackage{xcolor}
\usepackage[paperwidth=11cm,paperheight=8cm,margin=6mm]{geometry}
\pagestyle{empty}

\begin{document}
\setlength\parindent{0pt}\small

\textbf{（1）线型套线型：无法断行，溢出右边界}\par\smallskip
\CJKunderline{\CJKsout{虚室 hello 生白，吉祥 world 止止虚室生白，吉祥止止 虚室 hello 生白，吉祥 world 止止虚室生白，吉祥止止}}\par

\bigskip

\textbf{（2）线型套符号型：正常断行}\par\smallskip
\CJKunderline{\CJKunderdot{虚室 hello 生白，吉祥 world 止止虚室生白，吉祥止止 虚室 hello 生白，吉祥 world 止止虚室生白，吉祥止止}}\par

\bigskip

\textbf{（3）分段各用一个线型命令：正常断行}\par\smallskip
\CJKunderline{虚室 hello 生白，吉祥 world 止止虚室生白，吉祥止止} \CJKsout{虚室 hello 生白，吉祥 world 止止虚室生白，吉祥止止}\par

\end{document}
